\section{Fixed-node control}
\label{sec:theory}

This section analyzes the full-scan fixed-$B$ Stage~A feature test at a fixed
node and then the two Stage~B threshold tests for a fixed feature. The lemma
states the standard superuniformity of a single $+1$ Monte Carlo permutation
$p$-value. The corollaries apply the feature correction across the fixed-node
Stage~A tests and the threshold correction across the candidate thresholds, and
establish fixed-node control of the max-type rule; the proposition gives the
exact size of adaptive stopping. In the notation of \cref{sec:stageA-stageB},
this is the reference case where every feature in $F_t$ is tested with the same
fixed nodewise budget $B$. All probabilities below condition on $(X_t,U)$, where
$U$ collects label-independent auxiliary randomness other than the Monte Carlo
permutation draws. The assumptions of \cref{app:assumptions} fix the node and
its sample (A0.1), make the node responses exchangeable given $(X_t,U)$ under
the complete nodewise null (A0.2), fix the tested features and the budget as
functions of $(X_t,U)$ alone (A0.3), require exactly $B$ permutations with no
optional stopping (A0.4), and count ties conservatively (A0.5);
\cref{app:theory-scope} states the scope of each result.

The fixed-node results assume that the node sample is exchangeable in the sense
of A0.2, and bootstrap resampling does not satisfy that assumption: duplicated
rows carry the same response, so the node responses are not exchangeable given
the node covariates and the resampling indices.
\citet{Strobl2007BiasRFVariableImportance} show that sampling with replacement
biases variable selection in conditional inference forests, while subsampling
without replacement does not. The guarantees below therefore cover the node
tests of subsampled or unsampled fits.

\begin{lemma}[$+1$ permutation $p$-values are superuniform; \citealp{HemerikGoeman2018ExactTestingRandomPermutations}]
\label{lem:plusone-superuniform}
Assume A0.1 through A0.5 from \cref{app:assumptions}. Under the nodewise
complete permutation null at the fixed node $t$, let $T_{\mathrm{obs}}$ be the
observed statistic of one feature and $T_1,\dots,T_B$ its values under $B$
independent uniform random permutations of $Y_t$, and define the right-tail
$+1$ Monte Carlo $p$-value
%
\begin{equation}
  p_{\mathrm{MC}}^{\ge}
  := \frac{1 + \sum_{b=1}^B \ind\{T_b \ge T_{\mathrm{obs}}\}}{B+1}.
\end{equation}
%
Then for every $\alpha \in [0,1]$,
%
\begin{equation}
  \PP(p_{\mathrm{MC}}^{\ge} \le \alpha \mid X_t,U) \le \alpha.
\end{equation}
\end{lemma}

This is the standard exchangeable-rank result for Monte Carlo permutation tests
\citep{NorthCurtisSham2002EmpiricalPValuesMonteCarlo,Ernst2004PermutationMethodsExactInference,HemerikGoeman2018ExactTestingRandomPermutations}.
The proof derives the exchangeability it uses:
\cref{lem:identity-vs-random} shows that replacing $T_{\mathrm{obs}}$ by the
statistic of one more random permutation leaves the law of the $p$-value
unchanged, and the $B+1$ statistics so obtained are exchangeable, so the
inclusive rank of the first is superuniform
(\cref{app:permutation-exchangeability,app:proof-plusone}). The left-tail
version follows by applying the same result to $-T$.

\begin{corollary}[Stage~A complete null control at a fixed node]
\label{cor:stageA-global-null}
Assume A0.1 through A0.5 from \cref{app:assumptions}. Under the nodewise complete
null at $t$, let $F_t$ be nonempty and write $m_t = |F_t| \ge 1$. Let Stage~A
reject when
%
\begin{equation}
  \min_{j \in F_t} p_{t,j} \le \alpha_{\mathrm{sel}}/m_t.
\end{equation}
%
Then
%
\begin{equation}
  \PP\!\left(\exists j \in F_t : p_{t,j} \le \alpha_{\mathrm{sel}}/m_t \,\middle|\, X_t,U\right)
  \le \alpha_{\mathrm{sel}}.
\end{equation}
\end{corollary}

Combining \cref{lem:plusone-superuniform} with the Bonferroni union bound over
the $m_t$ Stage~A tests gives fixed-node Stage~A control under the complete
null. The tree-growing rule uses a strict comparison, so its rejection event is
a subset of the nonstrict event above. The result is local: it does not cover
partial-null familywise error control, $p$-values for selected features, or
error control for fitted trees and forests.

\begin{corollary}[Stage~B per-threshold rule at a fixed node]
\label{cor:stageB-bonferroni}
Assume A0.1 through A0.5 from \cref{app:assumptions} at a fixed node $t$, with the tested feature and the finite candidate set $C$, $K = |C| \ge 1$, fixed as functions of $(X_t,U)$ alone, and a fixed budget $B$ with no optional stopping. For $c \in C$ let $p_{t,c}$ be the left-tail $+1$ Monte Carlo $p$-value of the child impurity at threshold $c$, with ties counted with $\le$. Under the nodewise complete permutation null, conditional on $(X_t,U)$,
%
\begin{equation}
  \PP\!\left(\exists c \in C : p_{t,c} \le \alpha_{\mathrm{split}}/K \,\middle|\, X_t,U\right)
  \le \alpha_{\mathrm{split}}.
\end{equation}
\end{corollary}

Each $p_{t,c}$ is a $+1$ $p$-value of a statistic that is a fixed function of
the permuted response vector, so \cref{lem:plusone-superuniform}, applied in its
left-tail form, covers it. The union bound over the $K$ candidates then gives the
claim, as in \cref{cor:stageA-global-null} with $K$ in place of $m_t$; the bound
does not need the $K$ $p$-values to be independent, although they share the
permutations. Threshold scanning, which tests candidates in order of the
observed impurity and stops at the first rejection, leaves the rejection event
unchanged. The scope is weak familywise control over the $K$ candidates for a
feature fixed without reference to $Y_t$; it is not a post-selection guarantee
for a feature that Stage~A chose from the same responses. Its cost is the
budget: a $+1$ $p$-value cannot fall below $1/(B+1)$, so the rule can reject only
when $B \ge K/\alpha_{\mathrm{split}} - 1$.

\begin{corollary}[Stage~B max-type control at a fixed node]
\label{cor:stageB-maxt}
Assume A0.1 through A0.5 from \cref{app:assumptions} at a fixed node $t$, with the tested feature and the finite candidate set $C$ fixed as functions of $(X_t,U)$ alone, a fixed budget $B$ with no optional stopping, and left-tail ties counted with $\le$. Let $p^{\max}_t$ be the max-type Stage~B $p$-value of \cref{eq:maxt-stat} with column moments computed over all $B+1$ rows. Under the nodewise complete permutation null, conditional on $(X_t,U)$, $\PP(p^{\max}_t \le \alpha \mid X_t,U) \le \alpha$ for every $\alpha\in[0,1]$.
\end{corollary}

The argument reduces to \cref{lem:plusone-superuniform} in four steps.
(1)~Under the complete null, the $B+1$ rows of the impurity matrix
$(S_{b,c})$, computed from the observed response and from $B$ independent
uniform permutations of it, are exchangeable
\citep{HemerikGoeman2018ExactTestingRandomPermutations}. This is
\cref{lem:identity-vs-random} applied to the vector-valued statistic
$c \mapsto S_{b,c}$ in $\mathbb{R}^{K}$. That lemma holds for a statistic with
values in any measurable space: its proof establishes the joint law of the
permuted response vectors and applies the statistic only afterwards, as one
fixed measurable function of each vector, so it never uses that the statistic
is scalar. (2)~Column standardization uses the mean and standard deviation of
all $B+1$ entries of each column, which are symmetric functions of the rows, so
the standardized matrix has exchangeable rows as well; the exclusion of
candidates with an empty child or zero column variance is likewise a function of
the whole column. (3)~The row minimum is one fixed function applied to every
row, so $T_0,\dots,T_B$ are exchangeable. (4)~The inclusive-rank $p$-value of an
exchangeable tuple is superuniform by the proof of
\cref{lem:plusone-superuniform}, applied to $-T$ for the left tail.

The moments must include row $0$: estimating them from the permuted rows alone
singles out the observed row, breaks the symmetry the argument uses, and can be
anti-conservative. The guarantee is weak familywise control over the $K$
candidates, the probability of declaring any split under the complete nodewise
null, and its scope is that of \cref{cor:stageB-bonferroni}. It is a fixed-$B$
statement and does not cover the max-type rule with adaptive stopping, whose
column moments and chosen threshold move with the stopping time.

\subsection{Size of the default adaptive stopping rule}
\label{sec:theory-adaptive}

Adaptive stopping applies to one permutation test at level $\alpha$ with budget
$B$, confidence $c$, and batch size $s$. It draws permutations independently and
uniformly, with replacement, in batches of $s$. At each batch boundary $m$ that
is a multiple of $s$ with $m \ge \lceil 1/\alpha \rceil$, and at $m=B$, it stops
when the $\mathrm{Beta}(1+k,\,1+m-k)$ posterior places mass at least $c$ on
either side of $\alpha$, where $k$ counts permuted statistics at least as
extreme as the observed one. It then reports $(k+1)/(m+1)$, and the tree rejects
when the reported value is strictly below $\alpha$. Let $\rho(\alpha,c,s,B)$ be
the probability of rejection when the same rule is applied to a P\'olya urn
started from one ball of each colour. The library default is $c=0.95$ and
$s=32$.

\begin{proposition}[Size of the adaptive stopping rule, exact under A0.1 to A0.3 at fixed $B$]
\label{prop:adaptive-size}
Assume A0.1 through A0.3 at a fixed node for one permutation test whose
statistic is a fixed function of each permuted response vector (a Stage~A test
or a per-threshold Stage~B test), run with adaptive stopping at level $\alpha$,
confidence $c$, batch size $s$, and fixed budget $B$. Suppose that at every
reachable stopping state a significance stop reports a value strictly below
$\alpha$ and a futility stop reports a value at least $\alpha$. Under the
complete nodewise null, conditional on $(X_t,U)$, three statements hold.
(i)~Continuous mixing: if the per-draw probability that a permuted statistic is
at least as extreme as the observed one is uniform on $(0,1)$, the size equals
$\rho(\alpha,c,s,B)$ exactly. (ii)~Finite orbit: for a statistic without ties on
the $N = n_t!$ permutations, that probability is uniform on
$\{1/N, 2/N, \dots, 1\}$, and the size lies in $[\rho - 1/N,\ \rho]$.
(iii)~Ties: counting ties as exceedances can only lower the size, so with ties
the size is at most the value in (ii); for a discrete response, where
permutations tie, only this bound applies.
\end{proposition}

\emph{Reduction to a function of the exceedance probability.} Write $\theta$ for
the per-draw exceedance probability given the observed statistic and the node
data. Given $\theta$, the draws are independent, so the exceedance count after
$m$ draws is $\mathrm{Binomial}(m,\theta)$. The stop and report decisions are
deterministic functions of $(m,k)$ at the batch boundaries, so the rejection
probability $g(\theta)$ is a fixed function of $\theta$.

\emph{Monotonicity.} The function $g$ is nonincreasing. Flipping one indicator
from zero to one raises $k$ at every later boundary. A futility stop (posterior
mass at least $c$ above $\alpha$) can then only come earlier, and a significance
stop (mass at least $c$ below $\alpha$) only later. At a significance stop or at
the full budget, the reported value $(k+1)/(m+1)$ is at least as large on the
flipped path. The remaining step, that a significance stop reports a value
strictly below $\alpha$ and a futility stop a value at least $\alpha$, is the
hypothesis of the proposition.

\emph{From $g$ to statements (i) to (iii).} Under the complete null the observed
statistic is exchangeable with the $N$ orbit values, so without ties its rank
from the top is uniform on $\{1,\dots,N\}$ and $\theta$ is uniform on
$\{1/N,\dots,1\}$, which is (ii). The continuous idealization (i) replaces that
grid by the uniform distribution on $(0,1)$. Under it the binomial mixture is
the P\'olya urn started from one ball of each colour, whose next indicator is
one with probability $(k+1)/(m+2)$, and $\rho = \int_0^1 g$ by a finite
recursion over the stopping states. For nonincreasing $g$ the grid average
$N^{-1}\sum_{j=1}^{N} g(j/N)$ is at most $\int_0^1 g$ and at least
$\int_0^1 g - (g(0)-g(1))/N \ge \rho - 1/N$, which is the interval in (ii).
Counting ties as exceedances raises $\theta$, and given $\theta$ the same
nonincreasing $g$ gives (iii).

\begin{remark}[Values at the default settings]
An exact enumeration\footnote{The script
\texttt{paper/theory/adaptive\_stopping\_size.py} in the replication materials
regenerates the values of $\rho$, and the accompanying test verifies that the
implementation's Beta distribution function makes the same stopping decision as
a reference implementation at every reachable state.} verifies the hypothesis
on stopping reports at $c=0.95$ and $s=32$ for every budget from 20 to 3{,}000
at $\alpha=0.05$ and at every budget quoted in this remark. At $\alpha=0.05$,
$\rho(0.05,0.95,32,999)=0.0488$, $\rho(0.05,0.95,32,52)=0.0377$ at the auto
budget of 52 from \cref{eq:auto-budget}, and $\rho\le 0.0496$ for every $B$ from
20 to 3{,}000, a numerical statement whose largest value is at $B=140$; at
confidence 0.80, $\rho(0.05,0.80,32,999)=0.0522$. At the corrected levels the
tree tests, $\rho$ stays below the level. At $\alpha/20$, the feature correction
for 20 features, $\rho$ is $0.00249$ at the minimum budget of 400 and $0.00221$
at the auto budget of 3{,}144. At $\alpha/257$, the threshold correction for 257
candidates, $\rho$ is $0.000195$ at the minimum budget of 5{,}140 and $0.000181$
at the auto budget of 64{,}669. These values are for the strict rule the
implementation uses. Under the nonstrict convention, rejecting when the reported
value is at most $\alpha$, the corresponding values are $0.0495$ at $B=999$ and
$0.0523$ at $c=0.80$, and the bound $\rho \le \alpha$ fails at $B=79$, where
$\rho=0.05001$.
\end{remark}

\begin{remark}
Under the minimum budget with the feature or threshold correction, the first
stopping check at $\lceil 1/\alpha_{\mathrm{adj}}\rceil$
(\cref{sec:stageA-stageB}) coincides with the full budget, and the rule reduces
to the fixed-$B$ test of \cref{lem:plusone-superuniform}. This is the benchmark
configuration: every Stage~A and per-threshold Stage~B $p$-value in the
benchmark is a fixed-$B$ $+1$ permutation $p$-value. The max-type rule with
adaptive stopping recomputes its column moments over all permutations drawn so
far at each check, so its exceedance path is not a P\'olya urn and
\cref{prop:adaptive-size} does not cover it; its calibration rests on the null
simulations of \cref{sec:results-stageB}.
\end{remark}
